\icode{N23Y-MG-MU-L3-261021-001}
\nceeID[5]{2023重庆}{10}
（多选）
某卫星绕地心的运动视为匀速圆周运动，其周期为地球自转周期 \(T\) 的 \(\frac{3}{10}\)，运行的轨道与地球赤道不共面，如图所示。\(t_0\) 时刻，卫星恰好经过地球赤道上 \(P\) 点正上方。地球的质量为 \(M\)，半径为 \(R\)，引力常量为 \(G\)。则\paren[]

\begin{tikzpicture}[font=\scriptsize]
    \begin{scope}[rotate=45]
		\draw (-2.5,0) arc(180:0:2.5 and .6);
		\draw[fill=white] (0,0) circle(1.5);
		\draw (2.5,0) arc(0:-180:2.5 and .6);
	\end{scope}
	\draw[name path=c] (1.5,0) arc(0:-180:1.5 and .5);
	\draw[dashed] (-1.5,0) arc(180:0:1.5 and .5);
	\path[name path=l] (0,0)--++(-100:1);
	\path[name intersections={of=c and l,by=cd}];
	\fill (0,0) circle(.04) (cd)node[below left]{\(P\)} circle(.04);
\end{tikzpicture}
\begin{choices}
	\item 卫星距地面的高度为 \(\left(\frac{G M T^2}{4\pi^2}\right)^{\frac{1}{3}} - R\)
	\item 卫星与位于 \(P\) 点处物体的向心加速度大小比值为 \(\frac{5}{9\pi R} (180\pi G M T^2)^{\frac{1}{3}}\)
	\item 从 \(t_0\) 时刻到下一次卫星经过 \(P\) 点正上方时，卫星绕地心转过的角度为 \(20\pi\)
	\item 每次经最短时间实现卫星距 \(P\) 点最近到最远的行程，卫星绕地心转过的角度比地球的多 \(7\pi\)
\end{choices}
\begin{answer}
	BCD\qquad 提示：
\end{answer}